\section{Giuliano Imperatore}
\begin{quotex}
In the November 1940 issue of \emph{La Vita Italiana}, Julius Evola, writing under the pseudonym ``Arthos'', addressed the ancient concept of race in the article ``Sulle origini e sul doppio volto del razzismo''. We here are interested in the Emperor Julian, known as the Apostate. In this article, Evola aligns himself with Julian and states his opposition to the scientistic, genetic view of the origins of races, nations, and castes.
\end{quotex}

Flavio Claudio Giuliano\footnote{\url{http://it.wikipedia.org/wiki/Flavio\_Claudio\_Giuliano}} was the last of the pagan Roman emperors, the nephew of Constantine, who had established Christianity as the religion of the Empire. After several successful military campaigns against the Germans and Gauls, Julian became Emperor following the death of his cousin Constantius and resided in Constantinople, the seat of the Eastern part of the Empire. In the tradition of Marcus Aurelius, Julian was a philosopher-king. Although raised as a Christian, he abdicated under the influence of neo-platonic philosophy\footnote{\url{http://www.traditioninaction.org/Cultural/A014cpManualCivility\_Bearing.htm}}, which he studied in Athens (and was the classmate of two future saints: St. Basil the Great and St. Gregory). Apparently, he was initiated into the Eleusinian Mysteries and Mithras. As Emperor, he tried to govern on these three planks:

\begin{itemize}


\item Roman Tradition

\item Hellenic Culture

\item Neoplatonic philosophy
\end{itemize}


In short, we have a unique man: warrior, philosopher, emperor, initiate who supported the fundamental supporting principles of Western civilization. So what went wrong? Julian was not content to be non-Christian; he became anti-Christian. This led him to rule by force and violence, something Evola has warned against. He forced Christians out of state administrative duties and created martyrs, so any positive contributions he made were forever forgotten under the rubric of Julian the Apostate. A more cautious man may have noticed that his program was already in effect.

\begin{itemize}


\item \textbf{Roman Tradition}. He was concerned about the parallel hierarchy in the Church that mirrored the leadership hierarchy. Well, because of the hierarchy, the Catholic church was able to sustain Roman tradition and law well into the Middle Ages.

\item \textbf{Hellenic Culture}. He forbade the Christian schools from teaching Greek literature. This just isolated the Christians from Hellenic culture. Prior to that, Homer's epic poems were considered divinely inspired, Alexander was considered a saint, and Plato a revealer of divine Truth.

\item \textbf{Neoplatonic philosophy}. Christian theology was founded on Neoplatonism and should have been encouraged.
\end{itemize}


Although he himself did not believe in the existence of the Roman gods (he was a monotheist), he nevertheless tried to re-establish their cult. The problem is that once the noble lie is no longer believed, you cannot just will it back into existence. All that is now moot, in any case, since the so-called Reformation of the Northern peoples did much to destroy the triple foundation of Western civilization. So let's move on to Evola's points. Evola summarizes Julian's position:

\begin{quotex}
Against the Christian conception of the generation of men from a single couple, this emperor held to the doctrine of multiple creations, that is, of the diverse origins of men, unequal in race and disposition in the primordial times of the immortal gods. This idea reflects traditions very familiar to the ancient aryan peoples, which accounts for the differences not only of race, but also of castes and lineages or nations, \emph{stirpe} , by fundamentally distinct divine or demonic symbolic figures.
\end{quotex}


Thus, in his monotheism and polygenism, Julian is aligned with the theories of Herman Wirth\footnote{\url{https://gornahoor.net/library/UltimaThule.pdf}}. It is important to note here that Julian claims that races, nations, and even castes have their origin in spiritual differences, \emph{not} because of genetic factors. This is why Evola quotes him approvingly as one of the advocates of the Traditional doctrine of spiritual races. The modern mind can see origins only in term of time, in the past. So to them, the origins of man appear to be singular from the genetic perspective; even more absurdly, the genetic view would trace man's origin to the orangutan\footnote{\url{https://gornahoor.net/?p=1604}}.

Evola next addresses Thomas Campanella's City of the Sun\footnote{\url{https://gornahoor.net/?p=518}}. Although the City is based on traditional principles, such as rule by the philosopher-king and solarity, Evola notices the beginning of a more modern idea of race. For the Solarians, breeding is planned in the same way that horses or dogs are bred. This Evola criticizes as a rationalistic and mechanical approach. Evola then relates Campanella's ``solar race'' to the the Aryan-Hyperborean ``super-race'', and to the conception in the Middle Ages:

\begin{quotex}
the idea of race, united from the beginning to the idea of a public service, constituted one of the fundamental conditions of the nobility of the Middle Ages; a condition considered so natural, that no one thought of having to express it in words, much less in an ``myth'' or a theory. Let's keep in mind that, in later periods, ``science'', in the modern sense, appears always as the enemy of this true aristocracy, faithful to the idea of race up until the period of the Enlightenment, the precise point where the scientistic myth, in strict connexion Masonic agitations, helped it to dig its grave.
\end{quotex}



\flrightit{Posted on 2011-03-15 by Cologero}

\begin{center}* * *\end{center}

\begin{footnotesize}
\begin{sffamily}
\texttt{EXIT on 2011-03-16 at 09:50 said:}

Actually, Julian emphasized the differences in races and nations due to the national gods who were subordinate to the supreme god (thus, contrary to what you have stated, he did believe in the gods), but the actual existence of these differences are what really matters, and because of these differences he rejected Christian universalism which he considered racial treason and supported the ancestral traditions. Christianity, he said, was a wicked fiction. It mixes Jewish mentalities with misunderstood Aryan teachings, steals quotes from the mouths of gentiles, rewords them, then places it all into the mouth of a fictional messiah which is presented as historical fact. If you ask me Julian was not forceful enough in dealing with this false religion. This modern anti-power mentality is wholly foreign to true European character. It was Heraclitus who said that ``Strife is the father and king of all... it is necessary to understand that war is universal, and order is fighting, and everything comes about through fighting and necessity'' for ``the lightning bold steers all things.''

\hfill

\texttt{Cologero on 2011-03-16 at 12:49 said:}

``Belief'' is the word in question. He didn't accept a literal belief in the stories. To repeat Evola, which would be unnecessary to careful readers, races are the accounted for by ``fundamentally distinct divine or demonic symbolic figures''

So we see conjoined ``divine'' with ``symbolic figures'' ... the gods are ``symbolic figures'' which is all we implied. That, plus monotheism, makes Julian's ``religion'' something other than the traditional cult. Again, exactly what we claimed.

So all you show that Julian had his opinions, but was ineffective in his project ... again what we said. There are alternate ways to win the battle and Julian allowed his passion to get in the way of his intellect.

\hfill

\texttt{Albion on 2011-03-16 at 14:43 said:}

Evola's objection to materialistic orgins repeats almost verbatim Oswald Spengler's objection to Darwin, which he characterized as essentially ``English'' (by which I believe he meant materialistic). Likewise, Aristotle compares unfavorably in his book to Plato. The fundamental point seems to be that the origin of freedom lies in the invisible realm, and not the other way around. The material world can never superimpose itself in a one-to-one relationship, or else it would destroy that more real kingdom. This is a very helpful post, Cologero; it sheds a lot of light on why Rome collapsed.

\hfill

\texttt{Albion on 2011-03-16 at 16:09 said:}

Evola was working backwards, and so (in a sense) was Julian. Certainly, the gods do not look to us as they did to our ancestors. Nevertheless, they can remain real, more real perhaps, and differently real than before. Whatever Julian thought, Evola was most likely right to eschew some kind of fundamentalist re-paganizing worldview. St. Ambrose of Milan wrote, concerning Rome, ``it is no shame to pass on to better things''. Yet this could also be applied to ``Christianity''. It is always from the higher perspective that such superimposition is possible. I think this is what Evola was driving at.

\hfill

\texttt{Graham on 2011-03-16 at 16:27 said:}

The criticism of Julian is only over the policy of setting pagans and Christians against eachother, when under Constantine they had been synthesizing. Perhaps history affords a symbolic judgement of his efforts, in the division of the two `halves' (geographical and spiritual) of the Empire, and their subsequent fates.

\hfill

\texttt{Graham on 2011-03-16 at 16:50 said:}

Albion,

The unconditioned influence of `genius' or `daimon' would seem to be somewhat beyond Spengler's purview. For Evola, the origin is always potentially present; not so for Spengler, which explains his fatalism. This is a key difference between him and Tradition, which should be emphasized for the sake of understanding the latter.

\hfill

\texttt{John Bardis on 2011-03-18 at 05:52 said:}

As to ``believing'' in the gods, an interesting book might be Jean Richer's \_Sacred Geography of the Ancient Greeks\_ (1967, ET 1994). From the back cover:

``Richer found Plato's ideal city repeated around the most important oracular centers on ancient Greece. He shows how Plato's description was a later codification of a much earlier practice of dividing geography into twelve regions under the patronage of the gods of the zodiac. Several such twelve-part divisions of the Greek Territories are presented here.''

The center, of course, was the place of philosophy and monotheism.

As to the idea that each nation has its own angel--that comes from the Old Testament. Also note the matter of twelve tribes.

Another book that might be relevant is Anthony Damiani's \_Astronoesis\_ (2000). This brings astrology into Plotinus' system. In other words, the twelve signs, but especially the ``odd'' movements of the planets, explain the transition from the One to the many.

\hfill

\texttt{John Bardis on 2011-03-18 at 12:26 said:}

That reminds me of Hegel's view of the Monarch, here described by Zizek:

``Let us recall how the King--this exemplar of point de capiton, this individual who `quilts' the social edifice--was conceptualized by Hegel: the King is undoubtedly the point of the `suture' of social totality, the point whose intervention transforms a contingent collection of individuals into a rational totality--yet precisely as such, as the point which `sutures' Nature and Culture, as the point at which a cultural-symbolic function (that of being a king) immediately coincides with a natural determination (who will be king is determined by nature, by biological lineage), the King radically `desutures' all other subjects; makes them lose their roots in some preordained organic social body that would fix their place in society in advance and forces them to acquire their social status by means of hard labour. It is therefore not sufficient to define the King as the only immediate junction of Nature and Culture--the point is rather that this very gesture by means of which the King is posited as their `suture' de-sutures all other subjects, makes them lose their footing; throws them into a void where they must, so to speak, create themselves.''

\hfill

\texttt{Cologero on 2011-03-19 at 00:50 said:}

Interesting, since Evola was quite familiar with Hegel, though he never accepted Hegel's conception of a determinate unfolding of consciousness in history. However, I think I'm missing Zizek's point: this first part makes sense, but who are ``all other subjects''? Are they outsiders who don't fit into the rational totality? Is the point that with Julian, suddenly the Christians were ``desutured'' from the Empire?

\hfill

\texttt{John Bardis on 2011-03-19 at 09:21 said:}

Actually a King and an Emperor are two different things.

In Hegel's political philosophy he describes a very modern conception of the state. The one odd thing about it, though, for an American at least, is that he reserves a place in the modern state for a king. So the position of king would be the only hereditary office that he knew about. So the king is what he is by nature. Everyone else in the state has to become whatever it is he or she will become.

Like most Germans of that time, for Hegel ancient Greece was the ideal. It was an \_immediate\_ concrete universal. The Roman empire was an abstract universal. The modern state is a \_mediated\_ concrete universal.

So a king gives the people of a given state its identity. But an emperor destroys the identity of all the various people in the empire.

So, then, Hegel's view is \_very\_ different from that of the Traditionalists. He was very much opposed to the Romantic thinkers of his time who glorified India and the Catholic Church.

So, really, it isn't particularly appropriate to bring Hegel into things here. But on this one point, he seems to have been the last of the modern Western thinkers who had any appreciation at all, even if only very slightly, for a hereditary office.

\hfill

\end{sffamily}
\end{footnotesize}

